% Gerado por scripts/gerar_brasil.py. Relação integral de 73 presenças.
\section{Âmbito Brasil}
\begin{frame}{Microeletrônica no Brasil}
\duas{\begingroup\def\SemRotulos{1}\figura{\mapaBrasil{1,2,3,4,8,9}}\endgroup}{\vspace{5mm}\textcolor{blue_theme}{\Huge 54}\ organizações\par\medskip\textcolor{blue_theme}{\Huge 13}\ polos\par\medskip Projeto • ferramentas • integração • fabricação • pesquisa}
\creditoBrasil
\end{frame}
\note{Os números são da planilha local fornecida, com 54 organizações e 73 presenças em 13 polos. Diferem do resumo de outra versão da página do repositório. Um nome pode aparecer em vários polos. O mapa é próprio, com contorno geográfico simplificado e localizações esquemáticas agregadas. As próximas páginas preservam toda a relação; destaque exemplos em vez de ler cada nome. Centros e grupos de pesquisa são distinguidos de empresas. Créditos do levantamento: Matheus Grossi, Brazil Silicon Roadmap.}
\begin{frame}{Âmbito Brasil · Norte e Nordeste}
\begin{columns}[T,onlytextwidth]
\begin{column}{.35\textwidth}
\figura{\mapaBrasil{1,3,4}}
{\tiny $\bullet$ empresa\par $\circ$ centro/grupo}
\end{column}
\begin{column}{.63\textwidth}
\footnotesize\setlength{\columnsep}{3mm}\begin{multicols}{2}\raggedright
\textcolor{blue_theme}{\textbf{Manaus}}\par\smallskip
\noindent\parbox{\linewidth}{$\circ$ Instituto SENAI de Inovação em Microeletrônica}\par
\noindent\parbox{\linewidth}{$\bullet$ ADATA}\par
\noindent\parbox{\linewidth}{$\bullet$ Zilia}\par
\noindent\parbox{\linewidth}{$\bullet$ Stoneridge}\par
\noindent\parbox{\linewidth}{$\circ$ HBR}\par
\noindent\parbox{\linewidth}{$\circ$ Instituto de Pesquisas ELDORADO}\par
\noindent\parbox{\linewidth}{$\bullet$ HT Micron}\par
\noindent\parbox{\linewidth}{$\bullet$ CCBS / TERA}\par
\medskip
\textcolor{blue_theme}{\textbf{Campina Grande}}\par\smallskip
\noindent\parbox{\linewidth}{$\bullet$ HWiT}\par
\medskip
\textcolor{blue_theme}{\textbf{Salvador}}\par\smallskip
\noindent\parbox{\linewidth}{$\circ$ SENAI CIMATEC}\par
\noindent\parbox{\linewidth}{$\circ$ HBR}\par
\medskip
\end{multicols}
\end{column}
\end{columns}
\creditoBrasil
\end{frame}
\note{Relação da planilha PT-BR fornecida. Os polos abrangem cidades e proximidades ou áreas de interior; os pontos não indicam endereços individuais. As categorias e incertezas de atuação estão preservadas nos CSVs do projeto. Algumas organizações atuam em módulos, ferramentas ou eletrônica aplicada, sem confirmação de projeto local de CI/RTL. Os nomes estão aqui para consulta; escolha um ou dois exemplos para explicar a diversidade da cadeia.}
\begin{frame}{Âmbito Brasil · Centro-Oeste e Minas Gerais}
\begin{columns}[T,onlytextwidth]
\begin{column}{.35\textwidth}
\figura{\mapaBrasil{2,5,6}}
{\tiny $\bullet$ empresa\par $\circ$ centro/grupo}
\end{column}
\begin{column}{.63\textwidth}
\footnotesize\setlength{\columnsep}{3mm}\begin{multicols}{2}\raggedright
\textcolor{blue_theme}{\textbf{Brasília}}\par\smallskip
\noindent\parbox{\linewidth}{$\circ$ HBR}\par
\noindent\parbox{\linewidth}{$\circ$ Instituto de Pesquisas ELDORADO}\par
\medskip
\textcolor{blue_theme}{\textbf{Interior de MG}}\par\smallskip
\noindent\parbox{\linewidth}{$\circ$ LC-e}\par
\noindent\parbox{\linewidth}{$\circ$ Inatel}\par
\noindent\parbox{\linewidth}{$\bullet$ Modirum | GESPI}\par
\noindent\parbox{\linewidth}{$\bullet$ Marelli}\par
\noindent\parbox{\linewidth}{$\bullet$ Stellantis}\par
\noindent\parbox{\linewidth}{$\bullet$ Multilaser}\par
\medskip
\textcolor{blue_theme}{\textbf{Belo Horizonte}}\par\smallskip
\noindent\parbox{\linewidth}{$\bullet$ Cadence}\par
\noindent\parbox{\linewidth}{$\bullet$ RFIDO Design}\par
\medskip
\end{multicols}
\end{column}
\end{columns}
\creditoBrasil
\end{frame}
\note{Relação da planilha PT-BR fornecida. Os polos abrangem cidades e proximidades ou áreas de interior; os pontos não indicam endereços individuais. As categorias e incertezas de atuação estão preservadas nos CSVs do projeto. Algumas organizações atuam em módulos, ferramentas ou eletrônica aplicada, sem confirmação de projeto local de CI/RTL. Os nomes estão aqui para consulta; escolha um ou dois exemplos para explicar a diversidade da cadeia.}
\begin{frame}{Âmbito Brasil · São Paulo · polo de Campinas}
\begin{columns}[T,onlytextwidth]
\begin{column}{.35\textwidth}
\figura{\mapaBrasil{7}}
{\tiny $\bullet$ empresa\par $\circ$ centro/grupo}
\end{column}
\begin{column}{.63\textwidth}
\footnotesize\setlength{\columnsep}{3mm}\begin{multicols}{2}\raggedright
\textcolor{blue_theme}{\textbf{Campinas}}\par\smallskip
\noindent\parbox{\linewidth}{$\circ$ Von Braun Labs}\par
\noindent\parbox{\linewidth}{$\bullet$ ADATA}\par
\noindent\parbox{\linewidth}{$\circ$ HBR}\par
\noindent\parbox{\linewidth}{$\bullet$ HWiT}\par
\noindent\parbox{\linewidth}{$\bullet$ EnSilica}\par
\noindent\parbox{\linewidth}{$\circ$ CPQD}\par
\noindent\parbox{\linewidth}{$\bullet$ Aedvices}\par
\noindent\parbox{\linewidth}{$\circ$ Instituto de Pesquisas ELDORADO}\par
\noindent\parbox{\linewidth}{$\bullet$ Stoneridge}\par
\noindent\parbox{\linewidth}{$\bullet$ Valeo}\par
\noindent\parbox{\linewidth}{$\bullet$ Cadence}\par
\noindent\parbox{\linewidth}{$\circ$ CNPEM}\par
\noindent\parbox{\linewidth}{$\circ$ CTI Renato Archer}\par
\noindent\parbox{\linewidth}{$\bullet$ HCLTech}\par
\medskip
\end{multicols}
\end{column}
\end{columns}
\creditoBrasil
\end{frame}
\note{Relação da planilha PT-BR fornecida. Os polos abrangem cidades e proximidades ou áreas de interior; os pontos não indicam endereços individuais. As categorias e incertezas de atuação estão preservadas nos CSVs do projeto. Algumas organizações atuam em módulos, ferramentas ou eletrônica aplicada, sem confirmação de projeto local de CI/RTL. Os nomes estão aqui para consulta; escolha um ou dois exemplos para explicar a diversidade da cadeia.}
\begin{frame}{Âmbito Brasil · São Paulo · outros polos}
\begin{columns}[T,onlytextwidth]
\begin{column}{.35\textwidth}
\figura{\mapaBrasil{10,13}}
{\tiny $\bullet$ empresa\par $\circ$ centro/grupo}
\end{column}
\begin{column}{.63\textwidth}
\footnotesize\setlength{\columnsep}{3mm}\begin{multicols}{2}\raggedright
\textcolor{blue_theme}{\textbf{Interior de SP}}\par\smallskip
\noindent\parbox{\linewidth}{$\bullet$ Macnica DHW}\par
\noindent\parbox{\linewidth}{$\bullet$ WG3}\par
\noindent\parbox{\linewidth}{$\circ$ Chip Inventor}\par
\noindent\parbox{\linewidth}{$\bullet$ Beontag}\par
\noindent\parbox{\linewidth}{$\bullet$ SPICTEC}\par
\noindent\parbox{\linewidth}{$\circ$ LSI-TEC}\par
\noindent\parbox{\linewidth}{$\bullet$ Zilia}\par
\medskip
\textcolor{blue_theme}{\textbf{S. José dos Campos}}\par\smallskip
\noindent\parbox{\linewidth}{$\bullet$ Boeing}\par
\noindent\parbox{\linewidth}{$\bullet$ Embraer}\par
\noindent\parbox{\linewidth}{$\bullet$ Thales Group}\par
\noindent\parbox{\linewidth}{$\bullet$ SIATT}\par
\noindent\parbox{\linewidth}{$\bullet$ IACIT}\par
\noindent\parbox{\linewidth}{$\circ$ INPE}\par
\noindent\parbox{\linewidth}{$\bullet$ Modirum | GESPI}\par
\noindent\parbox{\linewidth}{$\bullet$ Tecsys do Brasil}\par
\medskip
\end{multicols}
\end{column}
\end{columns}
\creditoBrasil
\end{frame}
\note{Relação da planilha PT-BR fornecida. Os polos abrangem cidades e proximidades ou áreas de interior; os pontos não indicam endereços individuais. As categorias e incertezas de atuação estão preservadas nos CSVs do projeto. Algumas organizações atuam em módulos, ferramentas ou eletrônica aplicada, sem confirmação de projeto local de CI/RTL. Os nomes estão aqui para consulta; escolha um ou dois exemplos para explicar a diversidade da cadeia.}
\begin{frame}{Âmbito Brasil · Paraná e Santa Catarina}
\begin{columns}[T,onlytextwidth]
\begin{column}{.35\textwidth}
\figura{\mapaBrasil{11,8}}
{\tiny $\bullet$ empresa\par $\circ$ centro/grupo}
\end{column}
\begin{column}{.63\textwidth}
\footnotesize\setlength{\columnsep}{3mm}\begin{multicols}{2}\raggedright
\textcolor{blue_theme}{\textbf{Curitiba}}\par\smallskip
\noindent\parbox{\linewidth}{$\bullet$ HCLTech}\par
\noindent\parbox{\linewidth}{$\bullet$ Thales Group}\par
\noindent\parbox{\linewidth}{$\circ$ Lactec}\par
\medskip
\textcolor{blue_theme}{\textbf{Florianópolis}}\par\smallskip
\noindent\parbox{\linewidth}{$\bullet$ Macnica DHW}\par
\noindent\parbox{\linewidth}{$\bullet$ Chipus}\par
\noindent\parbox{\linewidth}{$\bullet$ Synopsys}\par
\noindent\parbox{\linewidth}{$\circ$ OSIRIS}\par
\noindent\parbox{\linewidth}{$\circ$ Silicium}\par
\medskip
\end{multicols}
\end{column}
\end{columns}
\creditoBrasil
\end{frame}
\note{Relação da planilha PT-BR fornecida. Os polos abrangem cidades e proximidades ou áreas de interior; os pontos não indicam endereços individuais. As categorias e incertezas de atuação estão preservadas nos CSVs do projeto. Algumas organizações atuam em módulos, ferramentas ou eletrônica aplicada, sem confirmação de projeto local de CI/RTL. Os nomes estão aqui para consulta; escolha um ou dois exemplos para explicar a diversidade da cadeia.}
\begin{frame}{Âmbito Brasil · Rio Grande do Sul}
\begin{columns}[T,onlytextwidth]
\begin{column}{.35\textwidth}
\figura{\mapaBrasil{9,12}}
{\tiny $\bullet$ empresa\par $\circ$ centro/grupo}
\end{column}
\begin{column}{.63\textwidth}
\footnotesize\setlength{\columnsep}{3mm}\begin{multicols}{2}\raggedright
\textcolor{blue_theme}{\textbf{Porto Alegre}}\par\smallskip
\noindent\parbox{\linewidth}{$\bullet$ Impinj}\par
\noindent\parbox{\linewidth}{$\bullet$ SILVACO}\par
\noindent\parbox{\linewidth}{$\bullet$ HCLTech}\par
\noindent\parbox{\linewidth}{$\bullet$ AEL Sistemas}\par
\noindent\parbox{\linewidth}{$\bullet$ Grupo Tellescom}\par
\noindent\parbox{\linewidth}{$\bullet$ Valeo}\par
\noindent\parbox{\linewidth}{$\circ$ Instituto de Pesquisas ELDORADO}\par
\noindent\parbox{\linewidth}{$\bullet$ EnSilica}\par
\noindent\parbox{\linewidth}{$\circ$ itt Chip — Unisinos}\par
\noindent\parbox{\linewidth}{$\bullet$ CEITEC}\par
\noindent\parbox{\linewidth}{$\circ$ Instituto SENAI (RS)}\par
\noindent\parbox{\linewidth}{$\bullet$ InPlace}\par
\noindent\parbox{\linewidth}{$\bullet$ HT Micron}\par
\medskip
\textcolor{blue_theme}{\textbf{Interior do RS}}\par\smallskip
\noindent\parbox{\linewidth}{$\bullet$ TUNA Semicondutores}\par
\noindent\parbox{\linewidth}{$\circ$ GMicro UFSM}\par
\medskip
\end{multicols}
\end{column}
\end{columns}
\creditoBrasil
\end{frame}
\note{Relação da planilha PT-BR fornecida. Os polos abrangem cidades e proximidades ou áreas de interior; os pontos não indicam endereços individuais. As categorias e incertezas de atuação estão preservadas nos CSVs do projeto. Algumas organizações atuam em módulos, ferramentas ou eletrônica aplicada, sem confirmação de projeto local de CI/RTL. Os nomes estão aqui para consulta; escolha um ou dois exemplos para explicar a diversidade da cadeia.}
